\lecture{14}
\title{Fast Fourier Transform}

\begin{document}

The explanation of this lecture has been altered
from the approach taken in 6.300 so as to
assume the knowledge of the Fast Fourier Transform
covered in \href{../../6.1220/notes-html/22-23.html}{Lectures 22 and 23}
of 6.1220.

\subsection{Recitation: DFT = Polynomial Sampling}

We first provide an alternative perspective of the DFT,
as suggested by 6.1220.

\begin{theorem}{DFT = Polynomial Sampling}
    Consider a DT signal $x[n]$ of length $N$ and its DFT $X[k]$.
    Furthermore, consider the polynomial $p(z) := \sum_{n = 0}^{N - 1} x[n] z^n$
    and a primitive $N^{\text{th}}$ root of unity $\omega := e^{-2\pi j/N}$.

    Then $\frac{1}{N} p(\omega^k) = X[k]$ for all $k \in \{0, 1, \dots, N - 1\}$;
    that is, the DFT effectively just evaluates a polynomial at $N$ roots of unity.
\end{theorem}

\begin{proof}
    Follows immediately from the DFT analysis equations from
    \href{12.html#lecture-discrete-fourier-transform}{Lecture 12}. ~ $\blacksquare$
\end{proof}

We may now apply reasoning shown in the 6.1220 notes to argue the following.

\begin{theorem}{DFT = Inverse DFT}
    If $X[k]$ is the DFT of $x[n]$,
    then the DFT of $X[k]$ is $x'[n] := \frac{1}{N} \cdot x[(-n)\%N]$.

    In other words, the inverse DFT of $X[k]$
    is the same as the DFT of $X[k]$,
    multiplied by $N$ and with circular time-reversal.
\end{theorem}

\begin{proof}
    Let's view the DFT as simply a map from
    vectors of length $N$ to vectors of length $N$.
    The previous theorem indicates that
    such a map is actually a linear transformation
    given by multiplication by the DFT matrix $W$.
    \[
    \begin{bmatrix}
        X[0] \\ X[1] \\ X[2] \\ \vdots \\ X[N - 1]
    \end{bmatrix}
    = \frac{1}{N} \begin{bmatrix}
        p(\omega^0) \\ p(\omega^1) \\ p(\omega^2) \\ \vdots \\ p(\omega^{N - 1})
    \end{bmatrix}
    = \frac{1}{N} \underbrace{
        \begin{bmatrix}
            1 & 1 & 1 & \cdots & 1 \\
            1 & \omega & \omega^2 & \cdots & \omega^{N - 1} \\
            1 & \omega^2 & \omega^4 & \cdots & \omega^{2(N - 1)} \\
            \vdots & \vdots & \vdots & \ddots & \vdots \\
            1 & \omega^{N - 1} & \omega^{2(N - 1)} & \cdots & \omega^{(N - 1)(N - 1)}
        \end{bmatrix}
    }_{W} \begin{bmatrix} x[0] \\ x[1] \\ x[2] \\ \vdots \\ x[N - 1] \end{bmatrix}.
    \]
    Thus, the DFT of $X[k]$ is just $\frac{1}{N^2} \cdot W^2 \cdot \mathbf{x}$,
    where $\mathbf{x}$ is the column vector associated with $x[n]$.
    \begin{quote}
        \textbf{Lemma.} We have $W^2 = NP$,
        where $P$ is the permutation matrix
        encoding the map $n \mapsto (-n) \% N$.

        \emph{Proof:} Recall that $W\overline{W} = NI$ from the
        \href{../../6.1220/notes-html/22-23.html}{Lectures 22 and 23}
        notes from 6.1220 (or just verify this via straightforward computation).
        The result $W^2 = NP$ then follows from noting easily
        that $\overline{W} = WP$. $\square$
    \end{quote}
    Thus, the DFT of $X[k]$ is
    $\frac{1}{N^2} \cdot W^2 \cdot \mathbf{x}
    = \frac{1}{N} \cdot P \cdot \mathbf{x}$,
    which indeed corresponds to the signal
    $x'[n] := \frac{1}{N} \cdot x[(-n)\%N]$. ~ $\blacksquare$
\end{proof}

Thus, any result we discover about the DFT
has a natural ``dual'' in the context of the inverse DFT.

\subsection{Lecture: The FFT Algorithm}

Refer to \href{../../6.1220/notes-html/22-23.html}{Lectures 22 and 23} for details.
The conclusion is that---via a Divide and Conquer algorithm---we may
compute the DFT of a DT signal of length $N$ in $O(N \log N)$ time.

Note that 6.1220 adopts the convention that
$p(\omega^k) = X[k]$, rather than the 6.300 convention that
$\frac{1}{N} p(\omega^k) = X[k]$,
so in the ``Combine Step'' section,
we multiply by an extra factor of $\frac{1}{2}$
that is otherwise omitted in the 6.1220 explanation.

\begin{verbatim}[python]
from cmath import exp, pi

def FFT(x):
    """Returns the DFT of x, an array whose length is a power of two."""

    # Base Case
    N = len(x)
    if N == 1:
        return x

    # Divide and Conquer
    even_x = x[::2]
    odd_x = x[1::2]
    FFT_even_x = FFT(even_x)
    FFT_odd_x = FFT(odd_x)

    # Combine Step
    FFT_x = []
    for k in range(N // 2):
        FFT_x.append((FFT_even_x[k] + exp(-2j * pi * k/N) * FFT_odd_x[k]) / 2)
    for k in range(N // 2):
        FFT_x.append((FFT_even_x[k] - exp(-2j * pi * k/N) * FFT_odd_x[k]) / 2)

    return FFT_x
\end{verbatim}

An implementation of the inverse DFT
would look very similar,
except with the exponents negated
and the factor of $\frac{1}{2}$ removed.

\end{document}